\documentclass[12pt]{article}
\usepackage[T1]{fontenc}
\usepackage[utf8]{inputenc}
\usepackage{lmodern}
\usepackage{textcomp}
\usepackage{enumitem}
\usepackage{geometry}
\geometry{margin=1in}
\begin{document}
\begin{center}
Answer Key - Biology - Graduate Level Assessment 17 \
\small (Answers are ordered exactly as in the question paper.)
\end{center}

\section*{Multiple Choice}
\begin{enumerate}[label=\arabic*.]
\item \textbf{Answer:} B
\\ \textit{Options:} A) The neurilemma sheath absorbs toxins at the injury site to facilitate regeneration.; B) The neurilemma sheath provides a channel for the axon to grow back to its former position.; C) The neurilemma sheath fuses severed axon ends together.; D) The neurilemma sheath repels inhibitory molecules that prevent nerve growth.; E) The neurilemma sheath acts as a source of neural stem cells for regeneration.
\\ \textit{Note:} The neurilemma sheath is critical for nerve regeneration because it forms a supportive channel that guides the regrowth of the severed axon back to its original target, facilitating functional recovery.
\item \textbf{Answer:} C
\\ \textit{Options:} A) Birds are called a society because each bird takes a turn leading the flock.; B) Birds in a flock are unrelated; C) A flock of birds can be considered an elementary form of society called a motion group.; D) The collective intelligence of a flock is referred to as a society.; E) Birds form a society to perform individual activities
\\ \textit{Note:} A flock of birds is referred to as a "society" because it exemplifies a basic social structure, known as a motion group, where individuals interact and coordinate their movements for survival and social cohesion.
\item \textbf{Answer:} C
\\ \textit{Options:} A) The suspensor is derived from the basal cell.; B) Cotyledons are derived from the apical cell.; C) Shoot apical meristem formation occurs after seed formation.; D) Precursors of all three plant tissue systems are formed during embryogenesis.
\\ \textit{Note:} The statement is correct because shoot apical meristem formation occurs during embryogenesis prior to seed formation, making it inaccurate to say it happens after seed formation.
\item \textbf{Answer:} C
\\ \textit{Options:} A) The organism is extinct in the wild and cannot be harvested for vaccine production; B) The organism that causes smallpox is too costly to use in vaccine production; C) The use of the actual virulent variola virus may be dangerous if too much is administered; D) Smallpox virus is confined to research laboratories; E) There is no specific treatment for smallpox
\\ \textit{Note:} The use of the actual virulent variola virus in smallpox vaccination poses significant risks of severe disease or complications if an excessive dose is administered, which is why safer alternatives, such as the vaccinia virus, are utilized instead.
\item \textbf{Answer:} A
\\ \textit{Options:} A) The gills of a fish and the lungs of a human; B) The leaves of an oak tree and the spines of a cactus; C) The eyes of a human and the eyes of a squid; D) The roots of a tree and the legs of an animal; E) The presence of a notochord in all chordate embryos
\\ \textit{Note:} The gills of a fish and the lungs of a human exemplify convergent evolution because they are functionally analogous structures that evolved independently in different species to serve the same purpose of respiration, illustrating how similar environmental pressures can lead to similar adaptations in unrelated organisms.
\end{enumerate}

\section*{True / False}
\begin{enumerate}[label=\arabic*.]
\setcounter{enumi}{5}
\item \textbf{Answer:} False
\\ \textit{Options:} A) True; B) False
\\ \textit{Note:} We found no evidence to support the use of AG routinely in IBS patients. Improvement of clinical response at 4-week follow-up may suggest a long-term effect of unknown mechanism, but could also be attributed to non-responder drop out. Gastrointestinal (GI) side effects may be a coincidence in this study, but irritation of GI tract by AG administration cannot be excluded.
\item \textbf{Answer:} False
\\ \textit{Options:} A) True; B) False
\\ \textit{Note:} Mothers' BMI highly correlate with children's BMI-z-scores. The degree of child's obesity increases mothers' concern and food restriction behavior. While mothers of obese children have a high prevalence of obesity, maternal obesity was found to have no significant influence on feeding behavior of obese school children.
\item \textbf{Answer:} True
\\ \textit{Options:} A) True; B) False
\\ \textit{Note:} The present survey showed a weak relationship between serum ACE and the frequency of SH, the clinical relevance of which is unclear. This limits the proposed role for serum ACE as an index of risk for SH.
\item \textbf{Answer:} True
\\ \textit{Options:} A) True; B) False
\\ \textit{Note:} General practitioners should consider using patients' first names more often, particularly with younger patients.
\item \textbf{Answer:} True
\\ \textit{Options:} A) True; B) False
\\ \textit{Note:} Our results give first evidence that lowered beta E during alcohol withdrawal may contribute to anxiety as a common disturbance during this state.
\end{enumerate}

\section*{Long Form Response}
\begin{enumerate}[label=\arabic*.]
\setcounter{enumi}{10}
\item \textbf{Answer:} Stem reptiles are crucial in the evolutionary lineage of modern reptiles as they represent a distinct group of ancestors that exclusively gave rise to contemporary reptiles without producing any other descendant lineages. This exclusivity highlights their role as a pivotal evolutionary node, marking the transition from early amniotes to the diverse clades we recognize today, such as squamates, turtles, and crocodilians. Their adaptations, including improved locomotion and a more efficient respiratory system, laid the groundwork for the ecological success of modern reptiles. Understanding stem reptiles allows researchers to trace the morphological and genetic developments that characterize present-day reptilian forms, underscoring their significance as a foundational lineage in the reptilian evolutionary tree.
\\ \textit{Note:} correct because it emphasizes the unique evolutionary significance of stem reptiles as exclusive ancestors that transitioned early amniotes into the diverse modern reptilian clades, underscoring their pivotal role in reptilian evolution.
\item \textbf{Answer:} The mapping of vegetation formations or biomes serves as a crucial tool for climatologists in understanding the climatic conditions within the biosphere. By analyzing the distribution of different vegetation types, researchers can infer regional climate characteristics such as temperature, precipitation patterns, and seasonal variations. For example, the presence of tropical rainforests indicates high humidity and consistent rainfall, whereas deserts signify arid conditions with minimal precipitation. Furthermore, the transition zones between biomes, such as temperate forests and grasslands, can reveal gradients in climate and serve as indicators for climate change impacts. Thus, vegetation mapping not only aids in the classification of ecosystems but also provides essential insights into the broader climatic dynamics of the Earth.
\\ \textit{Note:} correct because it highlights the relationship between vegetation formations and climatic conditions, demonstrating how the presence of specific biomes can indicate distinct climate characteristics such as humidity and precipitation patterns, which are essential for climatological analysis.
\end{enumerate}

\vfill
\noindent\textit{End of Answer Key}
\end{document}